\chapter{Parameter Identification} \label{paramIdent}
%
%
Here are some notes about parameter identification in piezoelectricity using CFS++. 
The analysistype is paramIdent. The following analysis section needs to be included in the xml-file

\begin{center}
\begin{minipage}{0.6\textwidth}
\begin{tabbing}
\quad\=\\ \kill
\verb!<paramIdent>!\\
\> \verb!<calcImpedanceCurve>! yes/no \verb!</calcImpedanceCurve>!\\
\> \verb!<calcMechDisplCurve>! yes/no \verb!</calcMechDisplCurve>!\\
\> \verb!<mechDisplAtNode> ! positive integer \verb! </mechDisplAtNode>!\\
\> \verb!<dofOfMechDispl>! positive integer \verb! </dofOfMechDispl>! \\
\> \verb!<numFreq>! positive integer \verb!</numFreq>! \\
\> \verb!<startFreq>! positive integer \verb!</startFreq>! \\
\> \verb!<stopFreq>! positive integer \verb!</stopFreq>! \\
\> \verb!<maxNumberInnerIterations>! positive integer \verb!</maxNumberInnerIterations>! \\
\> \verb!<maxNrOuterIterations>! positive integer \verb!</maxNrOuterIterations>! \\
\> \verb!<artDataNoise>! positive double \verb!</artDataNoise>! \\
\verb!</paramIdent>!\\
\end{tabbing}
\end{minipage}
\end{center}


\noindent Description:\\
The best way to check the quality of your parameters is, of course, a comparison of your simulated impedance curve (or mechanical displacement curve) with measurements. Setting \\ \verb!<calcImpedanceCurve>! yes \verb!</calcImpedanceCurve>! \\ (\verb!<calcMechDisplCurve>!yes\verb!</calcMechDisplCurve>!) leads to a calculation of such a curve in the interval $[startFreq,stopFreq]$ with {\em numFreq} evaluation points. The output will be a file called {\em imped.dat} ({\em mechDispl.dat}) which can be evaluated with MATLAB/GNUPlot, ... . The file {\em imped.dat} contains the computed {\it frequency, amplitude, phase, real(impedance), imag(impedance)} and the simulated {\it charge} at the surface of the ceramic.  The calculation will be done before and during the parameter identification process, so you can see wheather your calculation is successfull or not.  Note, parameter values are not taken from the material file {\em mat.dat or mat.xml}. You have to set them in the file {\em measuredData.dat}, see the following subsection.
\\
Further settings \\
\noindent \verb!<maxNrOuterIterations>! 100 - 200 \verb!</maxNrOuterIterations>! 
\verb!<maxNumberInnerIterations>! 30-100 \verb!</maxNumberInnerIterations>! 
\verb!<artDataNoise>! 0.01-1.0 \verb!</artDataNoise>!
are actually of no importance, since the methods specify them on their own.

\subsection{Additional input files}
The file  {\em measuredData.dat} actually handles most of the identification process. Important: Each line in this file needs to be closed by $'/'$ and all numbers are separated by $','$.\\ \\
\verb! #  denotes a comment line !\\
\verb! 1) 1.0E+06, 2.0E+06, 3.0E+06, ... / ! \\
Numbers after $1)$ are frequencies at which measurements will be read out of the additional input-file {\it mess.dat} containing frequencies, amplitudes and phases of the investigated specimen. \\
\noindent 
\verb! 4) 1.32000E+11, 1.150000E+11, 7.10000E+10, 7.30000E+10,  ... 5.84000E-09,/ ! \\
\verb! i) 0.069500E+11,0.08940E+11, -0.08080E+11,-0.08960E+11,  ... 0.0830E-08,/!\\
\verb! P) 1,  1,  1,  1,  1,  1,   1,  1,  1,   1,/ !\\
\verb! Q) 0,  1,  0,  0,  0,  0,  0 ,  1,  0,   1,/ !\\

\noindent The row $4)$ contains initial guesses for the real-valued material parameters\\ $(c_{11}, c_{33}, c_{12}, c_{13}, c_{44}, e_{15}, e_{31}, e_{33}, \epsilon_{11}, \epsilon_{33})$  followed by $i)$ the corresponding imaginary parts. The booleans after $P)$ and $Q)$ decide wheather the parameter should be updated during the identification or not. Here, all real parts will be fitted, from the imaginary parts only the three values $c_{33},e_{33}$ and $\epsilon_{33}$. Be aware while determining complex material parameters, that the materialtype is set to complex values in the xml-File, .i.e.\\ \noindent \verb!<materialDataType type="imagMaterialParameter"/>! \\
\noindent The choice of $M)$ is responsable which identification algorithm will be used, see Table \ref{TableNewtonMethods} for possible options.
\begin{table}
\begin{center}
\begin{tabular}{|c|c|c|}
\hline \hline {\bf M} & {\bf Method} & {\bf Comments}\\ 
\hline 1 &  Nonlinear - Landweber's iteration & robust, slow in convergence\\ 
\hline 4 & Newton-CG & just for academic purpose \\ 
\hline 5 & Newton-Landweber & partially robust, slow \\ 
\hline 6 & Newton-LandweberC (complex material parameter) & partially robust, slow\\ 
\hline 7 & Newton-$\nu$-Methods & robust, good convergence \\ 
\hline 8 & Newton-$\nu$-MethodsC (complex material parameter) & robust, good convergence\\ 
\hline 9-11 & Optimal experiment routines & just for academic purpose \\
\hline 12 & least-squares & robust, good convergence \\
\hline 
\end{tabular} 
\label{TableNewtonMethods}
\caption{Different identification methods and their corresponding number in the file {\em measuredData.dat}.}
\end{center}
\end{table} 

\noindent Different norms can be chosen by $N)$. The weighted norm $2$ should be the norm of choice.
The criteria, which physical quantity shall be considered during the identification is set behind $6)$. The options are logarithmic values of either the surface charge or the modulus of the impedance. ($1$ or $2$). Measuremets are listed in the file {\em mess.dat} in three columns, frequency, amplitude and phase.
\\
\subsection{Output files}
Several output files will be created during the run of the program:\\ \begin{itemize}
\item \noindent {\em parFinal.dat}, here, you find the computed parameters during the identification. Since they are in the same form as those in the file {\em measuredData.dat} you can simply copy them into this file and reuse them for a new fit, considering different frequencies, for example,\\
\item \noindent {\em parLog.dat}, records the norm of the residual and the development of the parameters,\\
\item \noindent {\em imped.dat}, impedance curve of actual parameters, \\
\item \noindent {\em mechDispl.dat}, mechanical deflection curve of actual parameters, \\
\end{itemize}
